\paragraph{Acknowledgements.}
We thank 
Alex Tamkin,
Colin Raffel,
Dan Jurafsky,
Deep Ganguli,
Douwe Kiela,
Emily Dinan,
Eric Horvitz,
Girish Sastry,
Iason Gabriel,
James Manyika,
Jason Wei,
Jie Tang,
Judy Shen,
Miles Brundage,
Neil Band,
Nelson Liu,
Opher Lieber,
Pang Wei Koh,
Stella Biderman,
Steven Cao,
Susan Zhang,
Teven Le Scao,
and
Yi Tay
for their valuable feedback on the manuscript.
We thank Steven Cao and Nelson Liu for guidance in the selection of \core scenarios that we describe in \refsec{core-scenarios}, and Kathy McKeown for her guidance on the summarization section.
We thank
Carlos Guestrin,
Daniel Zhang,
John Hewitt,
Kawin Ethayarajh,
Lauren Gillespie,
Mina Lee,
Rob Reich,
Rohan Taori,
Sandra Luksic,
Shelby Grossman, 
and
Yann Dubois
for helpful feedback and support throughout this project. 
We thank the CRFM community for helpful feedback on the effort overall. 

\paragraph{Model providers.}
We thank the following individuals at their respective organizations for the access, support, and/or credits required to evaluate their models:
\begin{itemize}
    \item AI21 Labs. Opher Lieber, Barak Lenz, Dan Padnos, Yoav Shoham
    \item Anthropic. Ben Mann, Jackson Kernion, Deep Ganguli, Jack Clark, Dario Amodei
    \item Cohere. Lewis Stott, Ellie Evans, Bill McCartney, Aidan Gomez
    \item Microsoft and the Turing Academic Program. Payal Bajaj, Barun Patra, Ahmed H. Awadallah, Saurabh Tiwary, Eric Horvitz
    \item OpenAI. Lama Ahmad, Miles Brundage
\end{itemize}
We thank CoreWeave for providing API access to \gptj and \gptneox for initial prototyping.

\paragraph{Funding and major compute.}
We thank Google through the Stanford HAI-Google collaboration for providing funding for this work, and especially Sebastian Gehrmann for overall feedback on this effort.
In addition, this work was supported in part by the AI2050 program at Schmidt Futures (Grant G-22-63429).
We thank Together Computer (which is backed by compute from Stanford University, ETH Zurich, Open Science Grid, University of Wisconsin-Madison, and Crusoe Cloud) for providing the infrastructure for benchmarking all the open models.
Rishi Bommasani was supported by the NSF Graduate Research Fellowship Program under grant number DGE-1655618.

\paragraph{Reflexivity.}
This work was developed at the Center for Research on Foundation Models (CRFM), a center at Stanford University borne out of the Stanford Institute for Human-Centered Artificial Intelligence (Stanford HAI).
It is accompanied by a policy brief on how evaluations substantively improve transparency.\footnote{See \url{https://hai.stanford.edu/foundation-model-issue-brief-series}.}

This work was made possible by two unique positions that CRFM occupies.
First, to develop the framework, build the benchmark, and execute the evaluation, CRFM serves as an interdisciplinary hub for coordinating and uniting the many authors of this work at Stanford University.
Second, to acquire the access and resources required to evaluate all the models in this work, CRFM leverages its relationships with the associated foundation model developers.

We further highlight that this work continues a trend documented by \citet{koch2021reduced} that many works that introduce (influential) benchmarks come from privileged, high-resourced, and powerful institutions, such as Stanford \citep[\eg][]{socher2013recursive, bowman2015large, rajpurkar2016squad, srivastava2021behavior}.
We actively encourage community-driven approaches to benchmark design, as well as mechanisms to center and drive adoption for benchmarks developed at other institutions.
This is of specific importance given how benchmarks embed values and priorities \citep{ethayarajh2020utility, birhane2022values}, which should be interrogated, as they are often developed by a small subset of the AI community but go on to orient the work of the broader AI community.
For this reason, we are very explicit about our design decisions (\eg prioritizing user-facing tasks in scenario selection) and directly foreground the limitations of our work with respect to the breadth of considerations in the AI community.
By fully releasing the open-source tools and raw model predictions, we invite everyone to contribute further scenarios, metrics, and models as we evolve \benchmarkname to better reflect the values of the AI community.